% theme: vectors_and_planar_motion
\MajorQuestionLesson{ベクトルと平面上の運動}
\SetRowSpace{1mm}

\TypeQuestionDouble{速度ベクトルの大きさと向きが次のように与えられている。$x$ 軸の正方向から反時計回りに測った角を用いて、各問いに答えなさい。}{vector_components}
\QQ{大きさが $12\mathrm{m/s}$、向きが $30^\circ$ である速度の $x$ 成分と $y$ 成分を求めなさい。}{$v_x=6\sqrt{3}\mathrm{m/s},\ v_y=6.0\mathrm{m/s}$}
\QQ{大きさが $10\mathrm{m/s}$、向きが $150^\circ$ である速度の $x$ 成分と $y$ 成分を求めなさい。}{$v_x=-5\sqrt{3}\mathrm{m/s},\ v_y=5.0\mathrm{m/s}$}
\QQ{成分が $(5\sqrt{3},-5.0)\mathrm{m/s}$ である速度の大きさと向きを求めなさい。}{$10\mathrm{m/s}$、$330^\circ$}
\QQ{成分が $(-4.0,4\sqrt{3})\mathrm{m/s}$ である速度の大きさと向きを求めなさい。}{$8.0\mathrm{m/s}$、$120^\circ$}

\TypeQuestionDouble{2つの速度ベクトルを $\vec{u}=(3.0,4.0)\mathrm{m/s}$、$\vec{v}=(-1.0,2.0)\mathrm{m/s}$ とする。}{vector_sum_difference}
\QQ{$\vec{u}+\vec{v}$ を成分で表しなさい。}{$(2.0,6.0)\mathrm{m/s}$}
\QQ{$\vec{u}+\vec{v}$ の大きさを求めなさい。}{$2\sqrt{10}\mathrm{m/s}$}
\QQ{$\vec{u}-\vec{v}$ を成分で表しなさい。}{$(4.0,2.0)\mathrm{m/s}$}
\QQ{$\vec{u}+\vec{w}=(0,5.0)\mathrm{m/s}$ となる速度ベクトル $\vec{w}$ を求めなさい。}{$\vec{w}=(-3.0,1.0)\mathrm{m/s}$}
\QQ{$2\vec{u}-\vec{v}$ を成分で表し、その大きさを求めなさい。}{$(7.0,6.0)\mathrm{m/s}$、$\sqrt{85}\mathrm{m/s}$}

\TypeQuestionDouble{平面上の点 ${\rm A}(1.0,2.0)$ から点 ${\rm B}(5.0,5.0)$ まで $2.0\mathrm{s}$ で移動し、続いて点 ${\rm C}(-1.0,8.0)$ まで $5.0\mathrm{s}$ で移動した。座標の単位は m とする。}{position_and_displacement}
\QQ{点 ${\rm A}$ から点 ${\rm B}$ までの変位ベクトルと、その大きさを求めなさい。}{$(4.0,3.0)\mathrm{m}$、$5.0\mathrm{m}$}
\QQ{点 ${\rm A}$ から点 ${\rm C}$ までの移動距離を求めなさい。}{$5.0+3\sqrt{5}\mathrm{m}$}
\QQ{点 ${\rm A}$ から点 ${\rm C}$ までの変位ベクトルと、その大きさを求めなさい。}{$(-2.0,6.0)\mathrm{m}$、$2\sqrt{10}\mathrm{m}$}
\QQ{全運動における平均の速度を成分で表しなさい。}{$\left(-\dfrac{2}{7},\dfrac{6}{7}\right)\mathrm{m/s}$}
\QQ{全運動における平均の速さを求めなさい。}{$\dfrac{5+3\sqrt{5}}{7}\mathrm{m/s}$}

\LessonPageBreak
\SetRowSpace{1mm}

\TypeQuestionDouble{ある物体が $t=0$ に点 $(-2.0,3.0)$ を通過し、その後、一定の速度 $\vec{v}=(4.0,-1.0)\mathrm{m/s}$ で運動する。位置の単位は m とする。}{uniform_planar_motion}
\QQ{時刻 $t$ における位置の $x$ 座標と $y$ 座標を、それぞれ $t$ の式で表しなさい。}{$x=-2.0+4.0t,\ y=3.0-t$}
\QQ{$2.0\mathrm{s}$ 後の物体の位置を求めなさい。}{$(6.0,1.0)\mathrm{m}$}
\QQ{物体が $x$ 軸上を通過する時刻と、そのときの位置を求めなさい。}{$3.0\mathrm{s}$ 後、$(10,0)\mathrm{m}$}
\QQ{物体が直線 $x=2y$ 上を通過する時刻と、そのときの位置を求めなさい。}{$\dfrac{4}{3}\mathrm{s}$ 後、$\left(\dfrac{10}{3},\dfrac{5}{3}\right)\mathrm{m}$}
\QQ{物体の速さを求めなさい。}{$\sqrt{17}\mathrm{m/s}$}

\TypeQuestionDouble{ある物体の時刻 $t$ における位置が、$x=t^2-4t+1$、$y=2t+3$ で表される。$t$ の単位は s、位置の単位は m とする。}{position_time_equations}
\QQ{$1.0\mathrm{s}$ 後の物体の位置を求めなさい。}{$(-2.0,5.0)\mathrm{m}$}
\QQ{速度の $x$ 成分と $y$ 成分を、それぞれ $t$ の式で表しなさい。}{$v_x=2t-4,\ v_y=2.0$}
\QQ{加速度の $x$ 成分と $y$ 成分を求めなさい。}{$a_x=2.0\mathrm{m/s^2},\ a_y=0\mathrm{m/s^2}$}
\QQ{$3.0\mathrm{s}$ 後の物体の速度と速さを求めなさい。}{$\vec{v}=(2.0,2.0)\mathrm{m/s}$、$2\sqrt{2}\mathrm{m/s}$}
\QQ{速度が $y$ 軸に平行になる時刻と、そのときの位置を求めなさい。}{$2.0\mathrm{s}$ 後、$(-3.0,7.0)\mathrm{m}$}
\QQ{$0\mathrm{s}$ から $2.0\mathrm{s}$ までの平均の速度を求めなさい。}{$(-2.0,2.0)\mathrm{m/s}$}

\TypeQuestionDouble{ある物体が $t=0$ に原点を速度 $\vec{v}_0=(6.0,8.0)\mathrm{m/s}$ で通過し、その後、一定の加速度 $\vec{a}=(-2.0,-2.0)\mathrm{m/s^2}$ で運動する。}{uniformly_accelerated_planar_motion}
\QQ{時刻 $t$ における速度と位置を成分で表しなさい。}{$\vec{v}=(6.0-2.0t,\ 8.0-2.0t)\mathrm{m/s}$、$\vec{r}=(6.0t-t^2,\ 8.0t-t^2)\mathrm{m}$}
\QQ{$2.0\mathrm{s}$ 後の速度と速さを求めなさい。}{$\vec{v}=(2.0,4.0)\mathrm{m/s}$、$2\sqrt{5}\mathrm{m/s}$}
\QQ{$2.0\mathrm{s}$ 後の物体の位置を求めなさい。}{$(8.0,12)\mathrm{m}$}
\QQ{速度が $x$ 軸に平行になる時刻と、そのときの位置を求めなさい。}{$4.0\mathrm{s}$ 後、$(8.0,16)\mathrm{m}$}
\QQ{物体の速さが最小になる時刻と、そのときの速さを求めなさい。}{$\dfrac{7}{2}\mathrm{s}$ 後、$\sqrt{2}\mathrm{m/s}$}
